\section{Tensor Products and the Tor Functor}
For the rest of the chapter, $R$ will denote a commutative ring, and we will work exclusively in $\Rmod{R} $.
\subsection{Bilinear Maps and the Definition of Tensor Product}
We have seen that $M\oplus N$ serves as both the product and coproduct of $R$-modules $M,N$, and thus an $R$-module homomorphism $M\oplus N\to P$ is determined by homomorphisms $M\to P$ and $N\to P$. However, there is another way to map the set $M\times N$ into $P$ in a way compatible with the $R$-module structure.
\begin{definition}[Bilinear]\label{dfn:8.8}
	Let $M,N,P$ be $R$-modules. A function $\varphi : M\times N \to P$ is $R$-bilinear if and only if
	\begin{itemize}
		\item For all $m\in M$, the map $n\mapsto \varphi (m,n)$ is an $R$-module homomorphism $N\to P$;
		\item for all $n\in N$, the map $m\mapsto \varphi (m,n)$ is an $R$-module homomorphism $M\to P$.
	\end{itemize}
\end{definition}

That is, if $\varphi $ is $R$-bilinear, then it is linear in both components. However, note that $\varphi $ itself is \emph{not linear}, even when viewed as $M\oplus N\to P$. However, we would expect there is a way to treat $\varphi $ as if it was $R$-linear. It is then natural to consider, using universal properties, a new $R$-module $M\otimes _RN$ together with an $R$-bilinear map
\[
	\otimes : M\times N \to M\otimes_R N
,\]
such that \emph{every} $R$-bilinear map factors uniquely through $M\otimes_R N$:
\[\begin{tikzcd}[row sep = 2.5em, column sep = 2.5em]
	M\times N\arrow[r, "\varphi "]\arrow[d, "\otimes "']&P\\
	M\otimes_R N\arrow[ru, dashed, "\overline{\varphi}"'] 
\end{tikzcd}\]
such that $\overline{\phi } $ is an $R$-module homomorphism. This is the tensor product.
\begin{definition}[Tensor Product]\label{dfn:8.9}
	The tensor product over $R$ of two $R$-modules $M,N$, denoted $M\otimes_R N$, is an $R$-module such that every $R$-bilinear map factors uniquely through it, in a way that the induced map is an $R$-linear map.
\end{definition}

The subscript $R$ is important, since $S$-bilinearity may be different from $R$-bilinearity, if $R\neq S$. It is not unusual to drop the subscript when the ring is understood, but we will generally not do this.

\begin{lemma}\label{lma:8.2}
	Tensor product exist in $\Rmod{R} $.
\end{lemma}
\begin{proof}
	Let $M,N$ be $R$-modules. Recall that the free $R$-module $F^R(M\times N)$ comes with a set map
	\[
		j : M\times N \to F^R(M\times N)
	\]
	which is universal with respect to functions $M\times N\to P$ for any $R$-module $P$. What we need is to make this into an $R$-bilinear map. This means we need to identify elements $j(m,n_1+n_2)$ with $j(m,n_1)+j(m,n_2)$, and similarly with $m$. Let $K$ be the free $R$-submodule generated by all elements of the form
	\[
		j(m,r_1n_1+r_2n_2)-r_2j(m,n_1)-r_2j(m,n_2)
	\]
	and
	\[
		j(r_1m_1+r_2m_2,n)-r_1j(m_1,n)-r_2j(m_2,n)
	.\]
	Define
	\[
		M\otimes_R N\coloneqq \frac{F^R(M\times N)}{K} 
	,\]
	and let $\otimes $ be the map $\pi j$, where $\pi $ is the natural projection:
	\[\begin{tikzcd}[row sep = 2.5em, column sep = 2.5em]
		\otimes :M\times N\arrow[r, "j"]&F^R(M\times N)\arrow[r, "\pi "]&\frac{F^R(M\times N)}{K} \eqqcolon M\otimes_R N
	\end{tikzcd}\]
	
	It is fairly obvious that this is an $R$-bilinear map at this point, since in the quotient module we will have
	\[
		j(m,r_1n_1+r_2n_2)-r_2j(m,n_1)-r_2j(m,n_2)=0\implies j(m,r_1n_1+r_2n_2)=r_1j(m,n_1)+r_2j(m,n_2)
	,\]
	and similarly for $m$. We now show that this construction satisfies the universal property. Let $\phi : M\times N \to P$ be any $R$-bilinear map. The universal property of free modules induces a homomorphism $\widetilde{\phi } $ such that
	\[\begin{tikzcd}[row sep = 2.5em, column sep = 2.5em]
		M\times N\arrow[d,"j"]\arrow[r, "\phi "]&P\\
		F^R(M\times N)\arrow[ru, dashed, "\widetilde{\phi }"]
	\end{tikzcd}\]
	commutes. Note that $\widetilde{\phi } $ will factor uniquely through the quotient by $K$ if $\ker \widetilde{\phi } =K$. By commutativity of the diagram, this is immediate:
	\begin{align*}
		\widetilde{\phi } &\left( j(m,r_1n_1+r_2n_2)-r_1j(m,n_1)-r_2j(m,n_2) \right) \\
				  &=\widetilde{\phi } \left( j(m,r_1n_1+r_2n_2 \right) -\widetilde{\phi } \left( r_1j(m,n_1) \right) -\widetilde{\phi } \left( r_2j(m,n_2) \right) \\
				  &=\phi (m,r_1n_1+r_2n_2)-r_1\phi (m,n_1)-r_2\phi (m,n_2)\\
				  &=0
	.\end{align*}
	The last step follows because $\phi $ is $R$-bilinear. It is easily shown that the converse is true; that if $\widetilde{\phi } (x)=0$, then $x$ is of this form. Therefore, $\widetilde{\phi } $ factors uniquely through the quotient by $K$:
	\[\begin{tikzcd}[row sep = 2.5em, column sep = 5em]
		M\times N\arrow[dd, bend right, in=-85, out=-75, "\otimes "']\arrow[d, "j"]\arrow[r, "\phi "]&P\\
		F^R(M\times N)\arrow[ru, "\widetilde{\phi }"]\arrow[d, "\pi "]\\
		M\otimes_R N\arrow[ruu, dashed, " \overline{\phi } "']
	\end{tikzcd}\]
	By commutativity of the diagram, we are done.
\end{proof}

One useful fact from this explicit construction is that elements of $M\otimes_R N$ arise from finite $R$-linear combinations
\[
	\sum_{i} r_i\left( m_i \otimes n_i \right) 
.\]
The $R$-bilinearity of the map $\otimes $ amounts to the statements that
\[
	m\otimes (n_1+n_2)=(m\otimes n_1)+(m\otimes n_2),
\]
\[
	(m_1+m_2)\otimes n=(m_1\otimes n)+(m_2\otimes n),
\]
\[
	m\otimes (rn)=(r m)\otimes n=r(m\otimes n)
.\]
It follows that elements of a tensor product can be viewed as
\[
	\sum_{i} r_i\left( m_i\otimes n_i \right) =\sum_{i} (r_im_i)\otimes n_i
.\]
Tensors that require only one term in the above summation; that is, tensors of the form $m\otimes n$, are called \emph{pure tensors}. Not all tensors are pure tensors, but pure tensors turn out to be somewhat useful, since they can generate the tensor product. For example, morphisms $\alpha ,\beta : M\otimes_R N \to P$ are the same if they coencide on all pure tensors.

\subsection{Adjunction with Hom and Explicit Computations}
Using the tensor product, we can define a collection of new covariant functors. Let $N$ be any $R$-module. Then the map $M\mapsto M\otimes_R N$ is a covariant functor $\Rmod{R} \to \Rmod{R} $. The action on morphisms is given as follows: for any morphism $\alpha : M_1 \to M_2$, there is an induced $R$-bilinear map $\alpha \times \mathrm{id}_N$, and thus an induced $R$-linear map
\[
	\alpha \otimes N : M_1 \otimes_R N \to M_2 \otimes_R N
.\]
This map is thus defined on pure tensors as $m\otimes n\mapsto \alpha (m)\otimes n$. Functorality follows immediately, and we have a covariant functor.

Recall that our definition of bilinear means that any $R$-bilinear map $\phi : M\times N \to P$ determines $R$-linear maps $\phi (-,n)$ and $\phi (m,-)$. That is, $\phi $ determines a map
\[
	M\to \Hom_{R}(N, P) 
.\]
Moreover, this map is an $R$-linear homomorphism (recall that hom-sets are $R$-modules with scalar multiplication given by $(r\phi )(x)=r\phi (x)$). That is, this bilinear map can be viewed as an element of
\[\Hom_{R}(M, \Hom_{R}(N, P) ) .\]
These considerations make the following statements fairly natural.
\begin{lemma}\label{lma:8.3}
	For all $R$-modules $M,N,P$, there is an isomorphism of $R$-modules
	\[
		\Hom_{R}(M, \Hom_{R}(N, P) ) \cong \Hom_{R}(M\otimes_RN, P) 
	.\]
\end{lemma}
\begin{proof}
	Every $\alpha \in \Hom_{R}(M, \Hom_{R}(N, P) ) $ determines an $R$-bilinear map by
	\[
		\phi : M\times N \to P,\qquad (m,n)\mapsto \alpha (m)(n)
	.\]
	By the universal property of tensors, $\phi $ factors uniquely through $\otimes $ as an $R$-linear map $\overline{\phi } : M\otimes_R N \to P$. It will be shown in the exercises that $\alpha \mapsto \overline{\phi } $ is linear, and an inverse will be constructed.
\end{proof}
\begin{corollary}\label{cor:8.1}
	For every $R$-module $N$, the functor $-\otimes_R N$ is left-adjoint to the functor $\Hom_{R}(N, -) $.
\end{corollary}
\begin{proof}
	By definition \ref{dfn:8.7}, these functors are adjoint if and only if
	\[
		\Hom_{R}(-\otimes_R N, \_) \cong \Hom_{R}(-, \Hom_{R}(N, \_) ) 
	.\]
	That is, for all $R$-modules $M,P$, we have
	\[
		\Hom_{R}(M\otimes_R N, P) \cong \Hom_{R}(M, \Hom_{R}(N, P) ) 
	.\]
	This is precisely the result of lemma \ref{lma:8.3}.
\end{proof}
It follows that the functor $-\otimes_R N$ preserves colimits, and by commutativity of the tensor product we have that $M\otimes_R -$ does as well. In particular, these functors are right-exact functors. These observations have many important consequences.
\begin{corollary}\label{cor:8.2}
	For all $R$-modules $M,N,P$,
	\[
		(M\oplus N)\otimes_R P \cong (M\otimes_R P)\oplus (N \otimes_R P)
	.\]
\end{corollary}
\begin{proof}
	Since $(M\oplus N)$ is the colimit of some functor, and since $-\otimes_R P$ is left-adjoint by corollary \ref{cor:8.1}, it will commute with coproducts by the dual of lemma \ref{lma:8.1}. In fact, it will commute for (possibly infinite) direct sums:
	\[
		\left( \bigoplus_{\alpha \in A} M_\alpha  \right) \otimes_R N \cong \bigoplus_{\alpha \in A}\left( M_\alpha \otimes_R N \right) 
	.\]
\end{proof}

This statement computes all tensors for tensor products of free $R$-modules.
\begin{corollary}\label{cor:8.3}
	For any two sets $A,B$,
	\[
		R^{\oplus A} \otimes_R R^{\oplus B} \cong R^{\oplus A\times B} 
	.\]
\end{corollary}
\begin{proof}
	By the previous corollary,
	\[
		R^{\oplus A} \otimes_R R^{\oplus B} \cong \left( R^{\oplus A}  \right) ^{\oplus B}
	.\]
	This has been shown previously to be isomorphic to $R^{\oplus A\times B} $.
\end{proof}
This shows that bases and generating sets work as expected - if $\left\{ e_i \right\}$ and $\left\{ f_i \right\} $ generate $M,N$ respectively, then the pure tensors $\left\{ e_i \otimes f_i \right\} $ generates $M\otimes_R N$. If the modules are free and these sets are minimal, then this is a basis for the tensor product. This is all that can really happen over a field - for finitely generates spaces of dimension $n$ and $m$, the tensor product just yields a space of dimension $mn$. Tensor products are thus more interesting over general rings.
\begin{corollary}\label{cor:8.4}
	For all $R$-modules $N$ and ideals $I$ of $R$,
	\[
		\frac{R}{I} \otimes_R N \cong \frac{N}{IN} 
	.\]
\end{corollary}
\begin{proof}
	Since $-\otimes_R N$ is right-exact, the sequence
	\[\begin{tikzcd}[row sep = 2.5em, column sep = 2.5em]
		0\arrow[r]&I\arrow[r]&R\arrow[r]&\frac{R}{I} \arrow[r]&0
	\end{tikzcd}\]
	induces the exact sequence
	\[\begin{tikzcd}[row sep = 2.5em, column sep = 2.5em]
		I\otimes_R N\arrow[r]&R\otimes_R N\arrow[r]&\frac{R}{I} \otimes_R N\arrow[r]&0
	\end{tikzcd}\]
	Note that $R\otimes _R N\cong N$. The image of $I\otimes_R N$ in this $R$-module is generated by the image of pure tensors $a\otimes n$ for $a\in I$, $n\in N$. This is $IN$ because $I$ is an ideal, and since the sequence is exact, the kernel of the $R$-module homomorphism $N\to \frac{R}{I} \otimes_R N$ is $IN$. This completes the proof by the first isomorphism thereom.
\end{proof}
\begin{corollary}\label{cor:8.5}
	For all ideals $I,J$ of $R$,
	\[
		\frac{R}{I} \otimes_R \frac{R}{J} \cong \frac{R}{I+J} 
	.\]
	
\end{corollary}
\begin{proof}
	This follows immediately from the previous corollary and the third isomorphism theorem.
\end{proof}

\subsection{Exactness Properties of Tensor; Flatness}
Although the tensor product is left-adjoint and thus right-exact, it is in general not an exact functor. For example, in the proof of corollary \ref{cor:8.4}, the map
\[
	I\otimes_R N\to N
\]
induced by the inclusion $I\hookrightarrow R$ may not be injective, and thus a short sequence would fail to be exact at $N$.

For another example, consider the injection $\mathbb{Z} \to \mathbb{Z} $ given by multiplication by 2. Tensoring by $\mathbb{Z} /2\mathbb{Z} $ over $\mathbb{Z} $, and keeping in mind that $R\otimes_R N\cong N$, we have the map (still given by multiplication by 2)
\[
	\frac{\mathbb{Z} }{2\mathbb{Z} } \to \frac{\mathbb{Z} }{2\mathbb{Z} } 
.\]
Note that this map sends $[0],[1]\mapsto 0$. This is thus the zero morphism, and clearly the exact sequence
\[\begin{tikzcd}[row sep = 2.5em, column sep = 2.5em]
	0\arrow[r]&\mathbb{Z} \arrow[r]&\mathbb{Z} \arrow[r]&\displaystyle{\frac{\mathbb{Z} }{2\mathbb{Z} }} \arrow[r]&0
\end{tikzcd}\]
is no longer exact after tensoring.

However, if $N$ is free, then $-\otimes_R N$ is exact. For every inclusion $M_1\hookrightarrow M_2$, tensoring by $N=R^{\oplus A} $ gives
\[\begin{tikzcd}[row sep = 2.5em, column sep = 2.5em]
	M_1^{\oplus A} \arrow[r, hook]&M_2^{\oplus A} 
\end{tikzcd}\]
which is given once again by the inclusion $M_1^{\oplus A} \subseteq M_2^{\oplus A} $. Due to this fact, tensoring is exact in $\Vect{k} $. That is, if
\[\begin{tikzcd}[row sep = 2.5em, column sep = 2.5em]
	0\arrow[r]&V_1\arrow[r]&V_2\arrow[r]&V_3\arrow[r]&0
\end{tikzcd}\]
is exact in $\Vect{k} $, then for all $k$-vector spaces $W$,
\[\begin{tikzcd}[row sep = 2.5em, column sep = 2.5em]
	0\arrow[r]&V_1\otimes_k W\arrow[r]&V_2\otimes_k W\arrow[r]&V_3\otimes_k W\arrow[r]&0
\end{tikzcd}\]
is exact. More simply, if $V_\bullet $ is exact, then $V_\bullet \otimes_k W$ is exact.
\begin{definition}[Flat]\label{dfn:8.10}
	An $R$-module $N$ is flat if the functor $\_\otimes_R N$ is exact.
\end{definition}
Flat modules are very important, especially in areas like algebraic geometry. It then makes sense to study the exactness of $R$-modules.

\subsection{The Tor Functors}
Now we will introduce the Tor functors, which are an important class of what is called a derived functor. The failiure of the functor $-\otimes_R N$ to be exact is measured by the other functor $\Tor_{1}^{R}(-, N) $. For example, if $N$ is flat, then $\Tor_{1}^{R}(M, N) =0$ for all $R$-modules $M$. Moreover, if
\[\begin{tikzcd}[row sep = 2.5em, column sep = 2.5em]
	0\arrow[r]&M\arrow[r]&N\arrow[r]&P\arrow[r]&0
\end{tikzcd}\]
is exact, then tensoring by any $A$ gives a new exact sequence
\[\begin{tikzcd}[row sep = 2.5em, column sep = 2.5em]
	\Tor_{1}^{R}(P, A) \arrow[r]&M\otimes_R A\arrow[r]&N\otimes_R A\arrow[r]&P\otimes_R A\arrow[r]&0
\end{tikzcd}\]
If the module on the left vanishes, then every exact sequence ending in $C$ remains exact under $-\otimes_R N$ in this case. This process can in fact be continued to obtain a longer exact complex
\[\begin{tikzcd}[row sep = 2.5em, column sep = 2em]
	\Tor_{1}^{R}(A, N) \arrow[r]&\Tor_{1}^{R}(B, N) \arrow[r]&\Tor_{1}^{R}(C, N) \arrow[r]&A\otimes_R N\arrow[r]&B\otimes _RN\arrow[r]&C\otimes _RN\arrow[r]&0
\end{tikzcd}\]
There in fact exist even more functors $\Tor_{2}^{R}(-, N) $, $\Tor_{3}^{R}(-, N) $ that continue this complex further. These are called the \emph{derived functors} of tensor.

This is how we compute these functors. Recall the definition of a free resolution of a module $M$: a free resolution of $M$ is an \emph{exact} complex $M_\bullet $ of free $R$-modules
\[\begin{tikzcd}[row sep = 2.5em, column sep = 2.5em]
	\cdots \arrow[r]&R^{\oplus S_2} \arrow[r]&R^{\oplus S_1} \arrow[r]&R^{\oplus S_0} \arrow[r]&M\arrow[r]&0.
\end{tikzcd}\]
If we throw away $M$ and tensor the free part by $N$, we get the complex $M_\bullet \otimes_R N$
\[\begin{tikzcd}[row sep = 2.5em, column sep = 2.5em]
	\cdots \arrow[r]&N^{\oplus S_2} \arrow[r]&N^{\oplus S_1} \arrow[r]&N^{\oplus S_0} \arrow[r]&0.
\end{tikzcd}\]
This follows from the fact that tensor commutes with colimits. The homology of this complex will suprisingly not depend (up to isomorphism) on the chosen free resolution, so we can define
\[
	\Tor_{i}^{R}(M, N) \coloneqq H_i(M_\bullet \otimes_R N)
.\]
By this definition, we will show in the exercises that $\Tor_{0}^{R}(M, N) =M\otimes _RN$, and if $N$ is flat, then $\Tor_{i}^{R}(M, N) =0$ for all $i>0$ (because tensoring by $N$ would be exact, so it would return an exact sequence, and would thus have no homology). In fact, this shows that if $\Tor_{1}^{R}(M, N) =0$ for all $M$, then $\Tor_{i}^{R}(M, N) $ for all $i$ and for all $M$. Then $N$ is flat.

We will later formalize these notions further, but for now we will simply assume definitions can be made as such. For now, we will do some simple computation with these functors. We have shown that every PID has a free resolution of length 1
\[\begin{tikzcd}[row sep = 2.5em, column sep = 2.5em]
	0\arrow[r]&R^{\oplus m_1} \arrow[r]&R^{\oplus m_0} \arrow[r]&M\arrow[r]&0
\end{tikzcd}\]
This property characterizes PIDs. Moreover, if
\[\begin{tikzcd}[row sep = 2.5em, column sep = 2.5em]
	0\arrow[r]&A\arrow[r]&B\arrow[r]&C\arrow[r]&0
\end{tikzcd}\]
is exact, then there exists compatible resolutions
\[\begin{tikzcd}[row sep = 2.5em, column sep = 2.5em]
	&0\arrow[d]&0\arrow[d]&0\arrow[d]\\
	0\arrow[r]&R^{\oplus a_1} \arrow[d, "\alpha "]\arrow[r]&R^{\oplus b_1} \arrow[d, "\beta "]\arrow[r]&R^{\oplus c_1} \arrow[r]\arrow[d, "\gamma "]&0\\
	0\arrow[r]&R^{\oplus a_0} \arrow[r]\arrow[d]&R^{\oplus b_0} \arrow[r]\arrow[d]&R^{\oplus c_0} \arrow[r]\arrow[d]&0\\
	0\arrow[r]&A\arrow[r]\arrow[d]&B\arrow[r]\arrow[d]&C\arrow[d]\arrow[r]&0\\
			   &0&0&0
\end{tikzcd}\]
such that columns and rows are exact. Tensoring the free rows by $N$ gives the exact diagram
\[\begin{tikzcd}[row sep = 2.5em, column sep = 2.5em]
	0\arrow[r]&N^{\oplus a_1} \arrow[r]\arrow[d, "\alpha \otimes N"]&N^{\oplus b_1} \arrow[r]\arrow[d, "\beta \otimes N"]&N^{\oplus c_1} \arrow[r]\arrow[d, "\gamma \otimes N"]&0\\
	0\arrow[r]&N^{\oplus a_0} \arrow[r]&N^{\oplus b_0} \arrow[r]&N^{\oplus c_0} \arrow[r]&0
\end{tikzcd}\]
Applying the snake lemma gives the exact sequence
\[\begin{tikzcd}[row sep = 2.5em, column sep = 1.5em]
	0\arrow[r]&H_1(A_\bullet \otimes_R N)\arrow[r]&H_1(B_\bullet \otimes _RN)\arrow[r]&H_1(C_\bullet \otimes _RN)\arrow[dlll, "\delta "]\\
	H_0(A\otimes _RN)\arrow[r]&H_0(B\otimes _RN)\arrow[r]&H_0(C_\bullet \otimes _RN)\arrow[r]&0
\end{tikzcd}\]
This is precisely the sequence of Tor modules
\[\begin{tikzcd}[row sep = 2.5em, column sep = 1.5em]
	0\arrow[r]&\Tor_{1}^{R}(A, N) \arrow[r]&\Tor_{1}^{R}(B, N) \arrow[r]&\Tor_{1}^{R}(C, N) \arrow[dlll, "\delta "]\\
	\Tor_{0}^{R}(A, N)\arrow[r]&\Tor_{0}^{R}(B, N) \arrow[r]&\Tor_{0}^{R}(C, N) \arrow[r]&0
\end{tikzcd}\]
The 0 on the left comes from the fact that $\Tor_{2}^{R}$ vanishes if $R$ is a PID. Moreover, $\Tor_{i}^{k}$ vanishes for all $i>0$ if $k$ is a field, since vector spaces are flat.
